%-----------PROJECTS-----------
\section{Projects}
    \vspace{-5pt}
    \resumeSubHeadingListStart
    \resumeProjectHeading
          {\textbf{\href{https://github.com/giorgosvyronos/hlp22-project}{ISSIE Extension}} $|$ \normalfont\textit{F\#, JavaScript, Electron}}{Jan 2022 --- Mar 2022}
          \resumeItemListStart
            \resumeItem{Developed in a group of 4, extended features for the university-led digital circuit design and simulation application, \uline{\href{https://tomcl.github.io/issie/}{ISSIE}}.}
            \resumeItem{Features included the ability for custom digital components to be drawn, moved, connected and rotated.}
            \resumeItem{Additional implementation included the design of a "properties" section, for renaming and remapping of component ports.}
        \resumeItemListEnd
            \vspace{-13pt}
      \resumeProjectHeading
          {\textbf{Custom FPGA SoC Design} $|$ \normalfont\textit{SystemVerilog, Intel Quartus Prime}}{Jan 2022 --- Mar 2022}
          \resumeItemListStart
            \resumeItem{Developed in a group of 2, a digital system design using Intel Quartus Prime to accelerate an application (arithmetic expression calculation) through reconfiguring the hardware of an FPGA.}
            \resumeItem{Designed a NIOS II processor on an FPGA using on-chip RAM and external memory.}
            \resumeItem{Extended design to provide support for floating point operations through software emulating blocks and through custom IP hardware blocks to be used as dedicated hardware blocks to compute the target expression.}
        \resumeItemListEnd
            % \vspace{-13pt}
    % \resumeProjectHeading
    %       {\textbf{IoT Microcontroller Game} $|$ \normalfont\textit{Eclipse, Intel Quartus Prime}}{Feb 2021 --- Mar 2021}
    %       \resumeItemListStart
    %         \resumeItem{Developed in a group of 6, a multiplayer racing game hosted on a server with multiple local FPGAs (DE10-lite) as nodes. The server was designed using NodeJS and Express.}
    %         \resumeItem{Steering data from the on-board accelerometer were locally processed and communicated to the server from each node. The data processing was designed using Eclipse. Server communication was established through Python and the Nios-II terminal.}
    %         \resumeItem{Information was then communicated back to each node to display current position on FPGA LEDs.}
    %       \resumeItemListEnd
        % \resumeProjectHeading
        %   {\textbf{IoT Microcontroller Game} $|$ \normalfont\textit{Eclipse, Intel Quartus Prime}}{Feb 2021 --- Mar 2021}
        %   \resumeItemListStart
        %     \resumeItem{Developed in a group of 6, a multiplayer racing game hosted on a server with multiple local FPGAs (DE10-lite) as nodes. The server was designed using NodeJS and Express.}
        %     \resumeItem{Steering data from the on-board accelerometer were locally processed and communicated to the server from each node. The data processing was designed using Eclipse. Server communication was established through Python and the Nios-II terminal.}
        %     \resumeItem{Information was then communicated back to each node to display current position on FPGA LEDs.}
        %   \resumeItemListEnd
    %        \vspace{-13pt}
    % \resumeProjectHeading
    %       {\textbf{MIPS-compatible CPU} $|$ \normalfont\textit{SystemVerilog, Intel Quartus Prime}}{Oct 2020 --- Nov 2020}
    %       \resumeItemListStart
    %         \resumeItem{Member of a group of 5 people that developed a working synthesizeable MIPS-compatible CPU using SystemVerilog. This CPU interfaced with the world using a memory-mapped bus, which would gave it access to memory and other peripherals.}
    %         \resumeItem{The outcome of this Project was to produce a deliverable containing an implementation of the CPU, a test-bench and memory, and a data-sheet covering the overall CPU architecture and the test-bench testing implementation.}
    %       \resumeItemListEnd
    \resumeSubHeadingListEnd
\vspace{-15pt}
